\documentclass[11pt,a4paper]{article}

% ── Packages ──────────────────────────────────────────────────────────────────
\usepackage[utf8]{inputenc}
\usepackage[T1]{fontenc}
\usepackage{lmodern}
\usepackage[margin=2.5cm]{geometry}
\usepackage{booktabs}
\usepackage{xcolor}
\usepackage{array}
\usepackage{multirow}
\usepackage{graphicx}
\usepackage{hyperref}
\usepackage{parskip}
\usepackage{fancyhdr}
\usepackage{titlesec}
\usepackage{amsmath}

% ── Colours ───────────────────────────────────────────────────────────────────
\definecolor{liu}{RGB}{0, 64, 140}        % LiU blue
\definecolor{oomred}{RGB}{200, 50, 50}
\definecolor{okgreen}{RGB}{34, 139, 34}
\definecolor{lightgray}{RGB}{245,245,245}

% ── Header / Footer ───────────────────────────────────────────────────────────
\pagestyle{fancy}
\fancyhf{}
\rhead{\textcolor{liu}{\small GPU Benchmark Report}}
\lhead{\small NVIDIA RTX PRO 6000 Blackwell}
\cfoot{\thepage}

% ── Section style ─────────────────────────────────────────────────────────────
\titleformat{\section}{\large\bfseries\color{liu}}{}{0em}{}[\titlerule]
\titleformat{\subsection}{\normalsize\bfseries}{}{0em}{}

% ── Helpers ───────────────────────────────────────────────────────────────────
\newcommand{\oom}{\textcolor{oomred}{\textbf{OOM}}}
\newcommand{\ok}{\textcolor{okgreen}{\checkmark}}

% ── Document ──────────────────────────────────────────────────────────────────
\begin{document}

% ── Title page ────────────────────────────────────────────────────────────────
\begin{titlepage}
  \centering
  \vspace*{2cm}
  {\Huge\bfseries\color{liu} GPU Performance Benchmark\\[0.4em]
   \large NVIDIA RTX PRO 6000 Blackwell Workstation\par}
  \vspace{1cm}
  {\large Linköping University --- AI Infrastructure Evaluation\par}
  \vspace{0.5cm}
  {\normalsize Date: March 2, 2026\par}
  \vspace{2cm}

  \begin{tabular}{ll}
    \toprule
    \textbf{GPU}          & NVIDIA RTX PRO 6000 Blackwell Workstation Edition \\
    \textbf{VRAM}         & 96 GB (97,887 MiB reported by driver) \\
    \textbf{Architecture} & Blackwell, Compute Capability 12.0 \\
    \textbf{System RAM}   & 134.97 GB \\
    \textbf{CPU Cores}    & 24 logical cores \\
    \textbf{TDP}          & 600 W \\
    \bottomrule
  \end{tabular}

  \vfill
  {\small Benchmarked using PyTorch 2.x, Hugging Face Transformers, and bitsandbytes.\\
  All results measured locally on the workstation with no network offloading.}
\end{titlepage}

% ── Abstract ──────────────────────────────────────────────────────────────────
\begin{abstract}
This report presents a comprehensive benchmark of the NVIDIA RTX PRO 6000 Blackwell
workstation GPU for Large Language Model (LLM) tasks. We evaluate three distinct
workloads: (1) raw hardware compute and memory bandwidth, (2) LLM inference throughput
across multiple model families and precisions, and (3) fine-tuning throughput under Full
Fine-Tuning, Low-Rank Adaptation (LoRA), and Quantized LoRA (QLoRA) regimes. System
resource utilisation (CPU, RAM, Swap) is monitored throughout. The results establish
effective capacity limits for both inference and training on a single-GPU workstation.
\end{abstract}

\tableofcontents
\newpage

% ══════════════════════════════════════════════════════════════════════════════
\section{Introduction}
% ══════════════════════════════════════════════════════════════════════════════

Modern Large Language Models range from a few hundred million to hundreds of billions of
parameters. Running or adapting them requires careful matching of model size to available
hardware. This report answers three practical questions:

\begin{enumerate}
  \item What is the raw compute capability of this GPU?
  \item Which LLMs can we run locally, and at what speed?
  \item Which LLMs can we fine-tune locally, and using which strategy?
\end{enumerate}

All benchmarks were run using open-source models downloaded from Hugging Face Hub.
Inference tests generated 200 tokens per run with a 27--32 token prompt.
Fine-tuning used the Stanford Alpaca dataset (52,000 instruction pairs) with
batch size 4 and sequence length 512.

% ══════════════════════════════════════════════════════════════════════════════
\section{Raw Hardware Compute}
% ══════════════════════════════════════════════════════════════════════════════

\subsection{Methodology}
Raw compute was measured by running repeated large matrix multiplications on an
$8192 \times 8192$ matrix — the dominant operation inside every transformer attention
and feedforward layer. Memory bandwidth was measured by copying a 2\,GB block between
GPU memory regions. A background CPU sampler recorded system load every 0.5\,s.

\subsection{Results}

\begin{table}[h!]
\centering
\caption{Raw GPU Compute and Memory Bandwidth}
\label{tab:raw}
\begin{tabular}{llrr}
\toprule
\textbf{Precision} & \textbf{Operation} & \textbf{Throughput} & \textbf{CPU Avg} \\
\midrule
FP32  & Matrix Multiply ($8192^2$) & 189.7 TFLOPS & 0.0\% \\
FP16  & Matrix Multiply ($8192^2$) & 300.9 TFLOPS & 13.5\% \\
BF16  & Matrix Multiply ($8192^2$) & 367.2 TFLOPS & 11.8\% \\
INT8  & Integer Multiply ($8192^2$) & 257.5 TOPS  & 11.7\% \\
\midrule
---   & Memory Bandwidth (2 GB copy) & \textbf{1{,}415 GB/s} & 12.0\% \\
\bottomrule
\end{tabular}
\end{table}

\subsection{Discussion}
BF16 is the preferred precision for LLM training due to its wider dynamic range
compared with FP16. The measured 367\,TFLOPS (BF16) significantly exceeds the
NVIDIA A100-80GB specification of 312\,TFLOPS. Memory bandwidth of 1{,}415\,GB/s
is critical because LLM \emph{decoding} (token generation) is memory-bound: the GPU
reads model weights for every token generated.

% ══════════════════════════════════════════════════════════════════════════════
\section{LLM Inference Benchmarks}
% ══════════════════════════════════════════════════════════════════════════════

\subsection{Methodology}
Each model was loaded once and ran a 200-token generation. We measure:
\begin{itemize}
  \item \textbf{TTFT} — Time To First Token (latency before output starts, ms).
  \item \textbf{Prefill tok/s} — Tokens per second while processing the input prompt (parallel).
  \item \textbf{Decode tok/s} — Tokens per second while generating new text (sequential, most important for user experience).
  \item \textbf{VRAM Peak} — Peak GPU memory allocated during inference.
\end{itemize}
Two precision levels are tested: \textbf{BF16} (full 16-bit weights) and
\textbf{4-bit NF4} (bitsandbytes NormalFloat quantisation).

\subsection{Results}

\begin{table}[h!]
\centering
\caption{LLM Inference Performance --- BF16 Precision}
\label{tab:inf-bf16}
\begin{tabular}{lrrrr}
\toprule
\textbf{Model} & \textbf{Params} & \textbf{TTFT (ms)} & \textbf{Decode tok/s} & \textbf{VRAM Peak} \\
\midrule
OPT-125M           & 0.1B  & 213  & \textbf{194.9} & 0.4 GB \\
OPT-1.3B           & 1.3B  & 12   & 103.3          & 2.9 GB \\
OPT-6.7B           & 6.7B  & 26   & 69.2           & 13.5 GB \\
OPT-13B            & 13B   & 24   & 49.7           & 25.9 GB \\
OPT-30B            & 30B   & 48   & 23.8           & 60.3 GB \\
\midrule
Qwen2.5-7B         & 7B    & 165  & 55.7           & 15.3 GB \\
Qwen2.5-14B        & 14B   & 41   & 33.4           & 29.6 GB \\
Qwen2.5-32B        & 32B   & 70   & 19.1           & 65.6 GB \\
Qwen2.5-Coder-7B   & 7B    & 59   & 57.0           & 15.3 GB \\
Qwen2.5-Coder-32B  & 32B   & 58   & 20.5           & 65.6 GB \\
\midrule
Mistral-7B         & 7B    & 34   & 56.2           & 14.5 GB \\
Mistral-Nemo-12B   & 12B   & 31   & 44.1           & 24.5 GB \\
Mistral-Small3-24B & 24B   & 50   & 29.5           & 47.2 GB \\
Mixtral-8x7B       & 46.7B & 1022 & 1.4 (slow)     & 87.3 GB \\
\midrule
DS-R1-Qwen-7B      & 7B    & 22   & 59.7           & 15.3 GB \\
DS-R1-Llama-8B     & 8B    & 51   & 55.4           & 16.1 GB \\
DS-R1-Qwen-14B     & 14B   & 52   & 35.1           & 29.6 GB \\
DS-R1-Qwen-32B     & 32B   & 59   & 20.5           & 65.6 GB \\
\bottomrule
\end{tabular}
\end{table}

\begin{table}[h!]
\centering
\caption{LLM Inference Performance --- 4-bit NF4 Quantization}
\label{tab:inf-4bit}
\begin{tabular}{lrrrr}
\toprule
\textbf{Model} & \textbf{Params} & \textbf{TTFT (ms)} & \textbf{Decode tok/s} & \textbf{VRAM Peak} \\
\midrule
OPT-125M           & 0.1B  & 156  & 77.4  & 0.2 GB \\
OPT-1.3B           & 1.3B  & 39   & 38.4  & 1.1 GB \\
Qwen2.5-7B         & 7B    & 61   & 27.1  & 5.7 GB \\
Qwen2.5-14B        & 14B   & 87   & 16.6  & 10.1 GB \\
Qwen2.5-32B        & 32B   & 133  & 12.5  & 19.6 GB \\
Qwen2.5-Coder-32B  & 32B   & 129  & 12.3  & 19.6 GB \\
Mistral-7B         & 7B    & 80   & 26.0  & 4.3 GB \\
Mistral-Nemo-12B   & 12B   & 72   & 20.4  & 8.5 GB \\
Mistral-Small3-24B & 24B   & 102  & 20.6  & 14.5 GB \\
DS-R1-Qwen-7B      & 7B    & 51   & 29.6  & 5.7 GB \\
DS-R1-Qwen-14B     & 14B   & 108  & 17.4  & 10.2 GB \\
DS-R1-Qwen-32B     & 32B   & 132  & 13.1  & 19.6 GB \\
\bottomrule
\end{tabular}
\end{table}

\newpage
\subsection{Discussion}
In BF16 mode, decoding speed decreases roughly as model size scales, since more
weight bytes must be read per token. Decode throughput above 30\,tok/s exceeds
the average adult reading speed (\textasciitilde{}250\,wpm $\approx$ 5\,tok/s readable),
making all 7B--14B models comfortable for interactive use. The 4-bit quantisation
reduces VRAM requirements by roughly $4\times$ at the cost of \textasciitilde{}40\%
lower throughput, which is the standard trade-off accepted in production deployments.

% ══════════════════════════════════════════════════════════════════════════════
\section{Fine-Tuning Benchmarks}
% ══════════════════════════════════════════════════════════════════════════════

\subsection{Methodology}
Fine-tuning training throughput measures how fast the GPU processes training samples
(forward pass + backward pass + optimiser step). Unlike inference, the GPU must
simultaneously hold weights, gradients, and optimiser states. We compared three regimes:

\begin{itemize}
  \item \textbf{Full Fine-Tuning (Full FT):} All $N$ parameters updated. Requires
        ${\sim}16\times$ the VRAM of the model weights alone (AdamW moments + gradients).
  \item \textbf{LoRA:} Only a small adapter ($<1\%$ of parameters) trained; base
        weights frozen in BF16. Reduces required VRAM by $3$--$4\times$.
  \item \textbf{QLoRA:} Base model loaded in 4-bit NF4; adapter trained in BF16.
        Further reduces VRAM by ${\sim}4\times$ compared with LoRA.
\end{itemize}

\subsection{Results}

\begin{table}[h!]
\centering
\caption{Fine-Tuning Throughput (Alpaca Dataset, BS=4, Seq=512)}
\label{tab:finetune}
\begin{tabular}{llrrr}
\toprule
\textbf{Model} & \textbf{Method} & \textbf{Tokens/s} & \textbf{VRAM Peak} & \textbf{Status} \\
\midrule
OPT-1.3B   & Full FT & 16{,}159 & 14.4 GB & \ok \\
OPT-1.3B   & LoRA    & 31{,}793 &  9.3 GB & \ok \\
OPT-1.3B   & QLoRA   & 13{,}783 &  3.6 GB & \ok \\
\midrule
Qwen2.5-7B & Full FT &  3{,}506 & 76.9 GB & \ok \\
Qwen2.5-7B & LoRA    &  7{,}200 & 36.8 GB & \ok \\
Qwen2.5-7B & QLoRA   &  3{,}299 & 14.6 GB & \ok \\
\midrule
Qwen2.5-14B & Full FT & ---     & \textgreater{}96 GB & \oom \\
Qwen2.5-14B & LoRA    & 3{,}964 & 64.2 GB & \ok \\
Qwen2.5-14B & QLoRA   & 1{,}890 & 21.9 GB & \ok \\
\midrule
DS-R1-Qwen-7B & Full FT & 3{,}504 & 76.9 GB & \ok \\
DS-R1-Qwen-7B & LoRA    & 7{,}187 & 36.8 GB & \ok \\
DS-R1-Qwen-7B & QLoRA   & 3{,}301 & 14.6 GB & \ok \\
\midrule
DS-R1-Qwen-14B & Full FT & --- & \textgreater{}96 GB & \oom \\
DS-R1-Qwen-14B & LoRA    & 3{,}960 & 64.2 GB & \ok \\
DS-R1-Qwen-14B & QLoRA   & 1{,}892 & 21.9 GB & \ok \\
\midrule
Phi-4-14B & Full FT & --- & \textgreater{}96 GB & \oom \\
Phi-4-14B & LoRA    & 4{,}520 & 54.0 GB & \ok \\
Phi-4-14B & QLoRA   & 2{,}274 & 17.1 GB & \ok \\
\bottomrule
\end{tabular}
\end{table}

\subsection{Discussion}
The OOM results for 14B Full FT are expected: AdamW maintains two momentum
buffers per parameter, so a 14B model ($\approx28$\,GB at BF16) requires
approximately $224$\,GB of VRAM just for weights + gradients + optimiser states.
No single GPU currently available can accommodate this. LoRA and QLoRA circumvent
this limitation effectively. Phi-4-14B achieved the highest LoRA throughput
(4{,}520 tok/s), attributed to its more efficient architecture design.

% ══════════════════════════════════════════════════════════════════════════════
\section{System Resource Utilisation}
% ══════════════════════════════════════════════════════════════════════════════

\subsection{Methodology}
A background thread sampled \texttt{psutil.cpu\_percent()} every 0.5\,s throughout
each benchmark stage. System RAM and Swap usage were captured at the end of each
stage. Low CPU utilisation indicates the GPU is not bottlenecked by data loading.

\subsection{Results}

\begin{table}[h!]
\centering
\caption{System Resource Utilisation During Benchmarks}
\label{tab:sysres}
\begin{tabular}{llrrr}
\toprule
\textbf{Task} & \textbf{Model Example} & \textbf{Avg CPU \%} & \textbf{RAM Used} & \textbf{Swap} \\
\midrule
Raw Compute (BF16 matmul) & ---              & 11.8\% & 24.7 GB & 0.0 GB \\
Raw Compute (INT8 matmul) & ---              & 11.7\% & 24.7 GB & 0.0 GB \\
Raw Bandwidth             & ---              & 12.0\% & 24.7 GB & 0.0 GB \\
\midrule
Inference                 & Qwen2.5-14B      &  11.4\% & 33.0 GB & 2.1 GB \\
Inference                 & DS-R1-Qwen-32B   &   9.6\% & 25.8 GB & 2.1 GB \\
Inference                 & Mistral-Small-24B &  9.5\% & 26.6 GB & 2.1 GB \\
\midrule
Fine-Tuning (Full FT)     & Qwen2.5-7B       &  8.5\% & 24.3 GB & 0.0 GB \\
Fine-Tuning (LoRA)        & DS-R1-Qwen-14B   &  9.8\% & 24.1 GB & 0.0 GB \\
Fine-Tuning (QLoRA)       & Phi-4-14B        &  9.2\% & 24.2 GB & 0.0 GB \\
\bottomrule
\end{tabular}
\end{table}

\subsection{Discussion}
CPU utilisation remained below 15\% in all workloads, confirming the GPU is the
sole bottleneck and the workstation is not CPU-starved. Swap usage of
$\approx2.1$\,GB during large inference runs is attributed to OS virtual memory
management for the model loader, not out-of-RAM conditions (134\,GB physical RAM
was never exhausted). Zero swap during fine-tuning confirms efficient GPU-side
batch construction.

% ══════════════════════════════════════════════════════════════════════════════
\section{Capacity Summary}
% ══════════════════════════════════════════════════════════════════════════════

\begin{table}[h!]
\centering
\caption{Single-GPU Capacity Limits (NVIDIA RTX PRO 6000, 96 GB)}
\label{tab:limits}
\begin{tabular}{llll}
\toprule
\textbf{Task} & \textbf{Precision} & \textbf{Max Model} & \textbf{Example} \\
\midrule
Inference & BF16 (16-bit)        & $\sim$40B params & Mixtral-8x7B (tight) \\
Inference & 4-bit NF4            & $\sim$140B params & Qwen-72B, LLaMA-70B \\
\midrule
Full Fine-Tuning & BF16          & $\sim$7–8B params & LLaMA-3.1-8B \\
LoRA Fine-Tuning & BF16          & $\sim$30–32B params & Qwen2.5-32B \\
QLoRA Fine-Tuning & 4-bit base  & $\sim$70B params & LLaMA-3.3-70B \\
\bottomrule
\end{tabular}
\end{table}

% ══════════════════════════════════════════════════════════════════════════════
\section{Conclusion}
% ══════════════════════════════════════════════════════════════════════════════

The NVIDIA RTX PRO 6000 Blackwell is a highly capable platform for local LLM
research and deployment. Key findings:

\begin{enumerate}
  \item \textbf{Compute:} 367\,TFLOPS (BF16) and 1{,}415\,GB/s bandwidth
        exceed the NVIDIA A100-80GB on both metrics.
  \item \textbf{Inference:} Models up to 32B parameters run at interactive
        speeds ($>$19\,tok/s) in full BF16 precision. 4-bit quantisation
        extends this to $\sim$140B without sacrificing usability.
  \item \textbf{Fine-Tuning:} Full fine-tuning is feasible up to 7--8B
        parameters. QLoRA enables adapter training on 14B--70B models with
        minimal quality degradation, making the workstation self-sufficient
        for most AI research workloads.
  \item \textbf{System Health:} CPU and RAM are not bottlenecks at any tested
        configuration. Swap usage is minimal and within expected OS behaviour.
\end{enumerate}

\end{document}
